%! TEX program = xelatex
\documentclass{beamer}
\usepackage{geometry}
\geometry{
  a4paper,
  total={170mm,257mm},
  left=20mm,
  top=20mm,
}
\usepackage{setspace}
\usepackage{graphicx}
\usepackage{xcolor}
\usepackage{kotex}
\usepackage{csquotes}
\usepackage{tabularray}
\usepackage{relsize}

%%% Typesetting for listings
\usepackage{listings}

\definecolor{dkgreen}{rgb}{0,0.6,0}
\definecolor{ltblue}{rgb}{0,0.4,0.4}
\definecolor{dkviolet}{rgb}{0.3,0,0.5}

%%% lstlisting coq style (inspired from a file of Assia Mahboubi)
\lstdefinelanguage{Coq}{ 
    % Anything betweeen $ becomes LaTeX math mode
    mathescape=true,
    % Comments may or not include Latex commands
    texcl=false, 
    % Vernacular commands
    morekeywords=[1]{Section, Module, End, Require, Import, Export,
        Variable, Variables, Parameter, Parameters, Axiom, Hypothesis,
        Hypotheses, Notation, Local, Tactic, Reserved, Scope, Open, Close,
        Bind, Delimit, Definition, Let, Ltac, Fixpoint, CoFixpoint, Add,
        Morphism, Relation, Implicit, Arguments, Unset, Contextual,
        Strict, Prenex, Implicits, Inductive, CoInductive, Record,
        Structure, Canonical, Coercion, Context, Class, Global, Instance,
        Program, Infix, Theorem, Lemma, Corollary, Proposition, Fact,
        Remark, Example, Proof, Goal, Save, Qed, Defined, Hint, Resolve,
        Rewrite, View, Search, Show, Print, Printing, All, Eval, Check,
        Projections, inside, outside, Def},
    % Gallina
    morekeywords=[2]{forall, exists, exists2, fun, fix, cofix, struct,
        match, with, end, as, in, return, let, if, is, then, else, for, of,
        nosimpl, when},
    % Sorts
    morekeywords=[3]{Type, Prop, Set, true, false, option},
    % Various tactics, some are std Coq subsumed by ssr, for the manual purpose
    morekeywords=[4]{pose, set, move, case, elim, apply, clear, hnf,
        intro, intros, generalize, rename, pattern, after, destruct,
        induction, using, refine, inversion, injection, rewrite, congr,
        unlock, compute, ring, field, fourier, replace, fold, unfold,
        change, cutrewrite, simpl, have, suff, wlog, suffices, without,
        loss, nat_norm, assert, cut, trivial, revert, bool_congr, nat_congr,
        symmetry, transitivity, auto, split, left, right, autorewrite},
    % Terminators
    morekeywords=[5]{by, done, exact, reflexivity, tauto, romega, omega,
        assumption, solve, contradiction, discriminate},
    % Control
    morekeywords=[6]{do, last, first, try, idtac, repeat},
    % Comments delimiters, we do turn this off for the manual
    morecomment=[s]{(*}{*)},
    % Spaces are not displayed as a special character
    showstringspaces=false,
    % String delimiters
    morestring=[b]",
    morestring=[d],
    % Size of tabulations
    tabsize=3,
    % Enables ASCII chars 128 to 255
    extendedchars=false,
    % Case sensitivity
    sensitive=true,
    % Automatic breaking of long lines
    breaklines=false,
    % Default style fors listings
    basicstyle=\small,
    % Position of captions is bottom
    captionpos=b,
    % flexible columns
    columns=[l]flexible,
    % Style for (listings') identifiers
    identifierstyle={\ttfamily\color{black}},
    % Style for declaration keywords
    keywordstyle=[1]{\ttfamily\color{dkviolet}},
    % Style for gallina keywords
    keywordstyle=[2]{\ttfamily\color{dkgreen}},
    % Style for sorts keywords
    keywordstyle=[3]{\ttfamily\color{ltblue}},
    % Style for tactics keywords
    keywordstyle=[4]{\ttfamily\color{dkblue}},
    % Style for terminators keywords
    keywordstyle=[5]{\ttfamily\color{dkred}},
    %Style for iterators
    %keywordstyle=[6]{\ttfamily\color{dkpink}},
    % Style for strings
    stringstyle=\ttfamily,
    % Style for comments
    commentstyle={\ttfamily\color{dkgreen}},
    %moredelim=**[is][\ttfamily\color{red}]{/&}{&/},
    literate=
    {\\forall}{{\color{dkgreen}{$\forall\;$}}}1
    {\\exists}{{$\exists\;$}}1
    {\ge -}{{$\geftarrow\;$}}1
    {=>}{{$\Rightarrow\;$}}1
    {==}{{\code{==}\;}}1
    {==>}{{\code{==>}\;}}1
    %    {:>}{{\code{:>}\;}}1
    {->}{{$\rightarrow\;$}}1
    {\ge ->}{{$\geftrightarrow\;$}}1
    {\ge ==}{{$\geq\;$}}1
    {\#}{{$^\star$}}1 
    {\\o}{{$\circ\;$}}1 
    {\@}{{$\cdot$}}1 
    {\/\\}{{$\text{ and }\;$}}1
    {\\\/}{{$\vee\;$}}1
    {++}{{\code{++}}}1
    {~}{{\ }}1
    {\@\@}{{$@$}}1
    {\\mapsto}{{$\mapsto\;$}}1
    {\\hline}{{\rule{\linewidth}{0.5pt}}}1
    %
}[keywords,comments,strings]

%%% Math settings
\usepackage{amssymb,amsmath,amsthm,mathtools}
\usepackage[math-style=TeX,bold-style=TeX]{unicode-math}
\theoremstyle{definition}
\newtheorem{definition}{Definition}[section]
\newtheorem{example}{Example}[section]
\newtheorem{lem}{Lemma}[section]
\newtheorem{thm}{Theorem}[section]
\newtheorem{cor}{Corollary}[section]
\newtheorem{clm}{Claim}[section]

%%% Font settings
%\setmainfont{Libertinus Serif}
%\setsansfont{Libertinus Sans}[Scale=MatchUppercase]
%\setmonofont{JuliaMono}[Scale=MatchLowercase]
%\setmainhangulfont{Noto Serif CJK KR}
%\setmonohangulfont{D2Coding}

%%% PL constructs
\usepackage{galois}
\usepackage{ebproof}
%% InfRule command
\ExplSyntaxOn
\NewTblrEnviron{@ruleenv}
\SetTblrInner[@ruleenv]{belowsep=0pt,stretch=0}
\SetTblrOuter[@ruleenv]{baseline=b}
\NewDocumentEnvironment { InfRule } { m +b }
  {
    \begin{@ruleenv}{l}
      \smaller{\textsc{#1}} \\
      \begin{prooftree} #2 \end{prooftree}
    \end{@ruleenv}
  }
  {}
\ExplSyntaxOff

%% Modification for \infer
\ExplSyntaxOn
\int_new:N \g__ebproof_sublevel_int
\box_new:N \g__ebproof_substack_box
\seq_new:N \g__ebproof_substack_seq

\cs_new:Nn \__ebproof_clear_substack:
  {
    \int_gset:Nn \g__ebproof_sublevel_int { 0 }
    \hbox_gset:Nn \g__ebproof_substack_box { }
    \seq_gclear:N \g__ebproof_substack_seq
  }

\cs_new:Nn \__ebproof_subpush:N
  {
    \int_gincr:N \g__ebproof_sublevel_int
    \hbox_gset:Nn \g__ebproof_substack_box
      { \hbox_unpack:N \g__ebproof_substack_box \box_use:c { \__ebproof_box:N #1 } }
    \seq_gput_left:Nv \g__ebproof_substack_seq
      { \__ebproof_marks:N #1 }
  }

\cs_new:Nn \__ebproof_subpop:N
  {
    \int_compare:nNnTF { \g__ebproof_sublevel_int } > { 0 }
      {
        \int_gdecr:N \g__ebproof_sublevel_int
        \hbox_gset:Nn \g__ebproof_substack_box {
          \hbox_unpack:N \g__ebproof_substack_box
          \box_gset_to_last:N \g_tmpa_box
        }
        \box_set_eq_drop:cN { \__ebproof_box:N #1 } \g_tmpa_box
        \seq_gpop_left:NN \g__ebproof_substack_seq \l_tmpa_tl
        \tl_set_eq:cN { \__ebproof_marks:N #1 } \l_tmpa_tl
      }
      { \PackageError{ebproof}{Missing~premiss~in~a~proof~tree}{} \__ebproof_clear:N #1 }
  }

\cs_new:Nn \__ebproof_append_subvertical:NN
  {
    \bool_if:NTF \l__ebproof_updown_bool
      { \__ebproof_append_above:NN #1 #2 }
      { \__ebproof_append_below:NN #1 #2 }
  }

\cs_new:Nn \__ebproof_join_subvertical:n
  {
    \group_begin:
    \__ebproof_subpop:N \l__ebproof_a_box
    \prg_replicate:nn { #1 - 1 }
      {
        \__ebproof_subpop:N \l__ebproof_b_box
        \__ebproof_enlarge_conclusion:NN \l__ebproof_b_box \l__ebproof_a_box

        \__ebproof_make_vertical:Nnnn \l__ebproof_c_box
          { \__ebproof_mark:Nn \l__ebproof_b_box {axis} - \__ebproof_mark:Nn \l__ebproof_b_box {left} }
          { \__ebproof_mark:Nn \l__ebproof_b_box {right} - \__ebproof_mark:Nn \l__ebproof_b_box {left} }
          { \skip_vertical:N \l__ebproof_rule_margin_dim }
        \__ebproof_vcenter:N \l__ebproof_b_box
        \__ebproof_append_subvertical:NN \l__ebproof_a_box \l__ebproof_c_box

        \__ebproof_append_subvertical:NN \l__ebproof_a_box \l__ebproof_b_box
      }
    \__ebproof_push:N \l__ebproof_a_box
    \group_end:
  }

\cs_new:Nn \__ebproof_renew_statement:nnn
  {
    \exp_args:Nc \RenewDocumentCommand { ebproof#1 }{ #2 } { #3 }
    \seq_gput_right:Nn \g__ebproof_statements_seq { #1 }
  }
\__ebproof_renew_statement:nnn { infer } { O{} m O{} m }
  {
    \group_begin:
    \__ebproof_restore_statements:
    \keys_set_known:nnN { ebproof / rule~style } { #1 } \l_tmpa_tl
    \keys_set:nV { ebproof } \l_tmpa_tl
    \tl_set:Nn \l__ebproof_right_label_tl { #3 }

    \__ebproof_clear_substack:
    \clist_map_inline:nn { #2 }
      {
        \__ebproof_join_horizontal:n { ##1 }
        \__ebproof_pop:N \l__ebproof_a_box
        \__ebproof_subpush:N \l__ebproof_a_box
      }
    \__ebproof_join_subvertical:n { \clist_count:n { #2 } }

    \__ebproof_push_statement:n { #4 }
    \__ebproof_join_vertical:
    \group_end:
  }
\ExplSyntaxOff
\ebproofset{left label template=\textsc{[\inserttext]}}
\ebproofset{center=false}

%%% Custom commands
\newcommand*{\vbar}{|}
\newcommand*{\finto}{\xrightarrow{\text{\textrm{fin}}}}
\newcommand*{\istype}{\mathrel{⩴}}
\newcommand*{\ortype}{\mathrel{|}}
\newcommand*{\cons}{::}
\newcommand*{\pset}{\mathscr{P}}
\newcommand*{\closeloc}[2]{\prescript{\text{\textbackslash}{#1}}{}{#2}}
\newcommand*{\openloc}[2]{{#1}^{#2}}

\def\ovbarw{1.2mu}
\def\ovbarh{1}
\newcommand*{\ovbar}[1]{\mkern \ovbarw\overline{\mkern-\ovbarw{\smash{#1}\scalebox{1}[\ovbarh]{\vphantom{i}}}\mkern-\ovbarw}\mkern \ovbarw}
\newcommand*{\A}[1]{\overset{\,_{\mbox{\Large .}}}{#1}}
\newcommand*{\Abs}[1]{{#1}^{\#}}
\newcommand*{\Expr}{\text{Expr}}
\newcommand*{\ExprVar}{\text{Var}}
\newcommand*{\Module}{\text{Module}}
\newcommand*{\ModVar}{\text{ModVar}}
\newcommand*{\modid}{d}
\newcommand*{\Time}{\mathbb{T}}
\newcommand*{\ATime}{\A{\Time}}
\newcommand*{\Ctx}{\text{Env}}
\newcommand*{\ctx}{\sigma}
\newcommand*{\Value}{\text{Val}}
\newcommand*{\Walue}{\text{WVal}}
\newcommand*{\Mem}{\text{Mem}}
\newcommand*{\Loc}{\text{Loc}}
\newcommand*{\FLoc}{\text{FLoc}}
\newcommand*{\Left}{\text{Left}}
\newcommand*{\Right}{\text{Right}}
\newcommand*{\Sig}{\text{Sig}}
\newcommand*{\mem}{m}
\newcommand*{\dom}{\text{dom}}
\newcommand*{\AMem}{\A{\text{Mem}}}
\newcommand*{\State}{\text{State}}
\newcommand*{\AState}{\A{\text{State}}}
\newcommand*{\Result}{\text{Result}}
\newcommand*{\AResult}{\A{\text{Result}}}
\newcommand*{\Tick}{\text{Tick}}
\newcommand*{\lfp}{\mathsf{lfp}}
\newcommand*{\Step}{\mathsf{Step}}
\newcommand*{\semarrow}{\Downarrow}
\newcommand*{\asemarrow}{\A{\rightsquigarrow}}
\newcommand*{\synlink}{\rtimes}
\newcommand*{\semlink}{\mathbin{\rotatebox[origin=c]{180}{$\propto$}}}
\newcommand*{\link}[2]{{#1}\rtimes{#2}}
\newcommand*{\mt}{\mathsf{empty}}
\newcommand*{\valid}{\mathsf{valid}}
\newcommand*{\Path}{\text{Path}}
\newcommand*{\equivalent}{\sim}

\newcommand*{\doubleplus}{\ensuremath{\mathbin{+\mkern-3mu+}}}
\newcommand*{\project}{\text{\texttt{:>} }}
\newcommand*{\Exp}{\mathsf{Exp}}
\newcommand*{\Imp}{\mathsf{Imp}}
\newcommand*{\Fin}{\mathsf{Fin}}
\newcommand*{\Link}{\mathsf{Link}}
\newcommand*{\sembracket}[1]{\lBrack{#1}\rBrack}
\newcommand*{\fin}[2]{{#1}\xrightarrow{\text{fin}}{#2}}
\newcommand*{\addr}{\mathsf{addr}}
\newcommand*{\tick}{\mathsf{tick}}
\newcommand*{\modctx}{\mathsf{ctx}}
\newcommand*{\mapinject}[2]{{#2}[{#1}]}
\newcommand*{\inject}[2]{{#2}\langle{#1}\rangle}
\newcommand*{\deletepre}[2]{{#2}\overline{\doubleplus}{#1}}
\newcommand*{\deletemap}[2]{{#1}\overline{[{#2}]}}
\newcommand*{\delete}[2]{{#2}{\langle{#1}\rangle}^{-1}}
\newcommand*{\filter}{\mathsf{filter}}
\newcommand*{\Lete}{\mathtt{val}}
\newcommand*{\Letm}{\mathtt{mod}}

\newcommand*{\ValRel}[1]{\mathcal{V}\sembracket{#1}}
\newcommand*{\ExprRel}[1]{\mathcal{E}\sembracket{#1}}
\newcommand*{\CtxRel}[1]{\mathcal{C}\sembracket{#1}}
\newcommand*{\ModRel}[1]{\mathcal{M}\sembracket{#1}}
\newcommand*{\TyEnv}{\text{TyEnv}}
\newcommand*{\TyVar}{\text{TyVar}}
\newcommand*{\Type}{\text{Type}}
\newcommand*{\Subst}{\text{Subst}}
\newcommand*{\external}{\Gamma_{\text{ext}}}

\newcommand*{\InitE}{\textsf{Init}}
\newcommand*{\ReadE}{\textsf{Read}}
\newcommand*{\CallE}{\textsf{Call}}
\newcommand*{\resolve}{\downarrow_{\ctx_0}}


\title{프로그램 따로분석의 이론적 기틀}
%\subtitle{}
\author{이준협}
\date{2024년 4월 19일}
\institute{ROPAS S\&T}

\titlegraphic{%
  \begin{tikzpicture}[overlay,remember picture]
    \node at (current page.145) [xshift=3em, yshift=-1.3em] {
      \includegraphics[width=3em]{snu-symbol.png}
    };
    \node at (current page.35) [xshift=-3em, yshift=-1.3em] {
      \includegraphics[width=2.5em]{ropas-symbol.png}
    };
  \end{tikzpicture}%
}

\begin{document}
\maketitle
\begin{frame}[c,fragile]
	\frametitle{따로분석?}
	\begin{itemize}
		\item 따로분석이란,
		\item 모든 것이 알려지기 전에,
		\item 프로그램의 의미를 포섭하는 분석이다.
	\end{itemize}
\end{frame}
\begin{frame}[c,fragile]
	\frametitle{왜 필요한가?}
	\begin{tabular}{c}
		\begin{lstlisting}[language=Simple]
int *x;
void f(void) {
  if (*x == *x) g(); /* call to unknown function */
  *x = 42; /* what happens here? */
}
  \end{lstlisting}
	\end{tabular}
	\begin{itemize}
		\item 코드 전체가 주어지지 않는 일이 빈번하다.\\
		      \textbf{ex. }외부 라이브러리 함수를 부르는 경우
		\item \textbf{미리} 최대한 분석하고 싶다.
	\end{itemize}
\end{frame}
\begin{frame}[c,fragile]
	\frametitle{미리보기}
	\begin{center}
		$\sembracket{e}(\Sigma_0\semlink\Sigma)\subseteq\Sigma_0\semlink\sembracket{e}\Sigma$
	\end{center}
	\begin{itemize}
		\item $\Sigma_0:\text{내가 모르던 외부 환경}$
		\item $\sembracket{e}\Sigma:\text{몰라도 최선을 다해 실행한 결과}$
		\item $\semlink:\text{합치기(linking) 연산}$
	\end{itemize}
	\[\text{따로분석}=\text{안전한 }\Abs{\sembracket{\cdot}}+\text{안전한 }\Abs\semlink\]
	\[\text{분석의 안전성}:\sembracket{e}(\Sigma_0\semlink\Sigma)\subseteq\gamma(\overbracket{\Abs\ctx_0\Abs\semlink\underbracket{\Abs{\sembracket{e}}\Abs\ctx}_{\text{따로분석}}}^{\text{나중에 합치기}})\]
\end{frame}
\begin{frame}[c,fragile]
	\frametitle{언어 소개}
	\begin{figure}
		\small
		\begin{tabular}{rrcll}
			Identifiers & $x$ & $\in$         & $\ExprVar$                                                      \\
			Expression  & $e$ & $\rightarrow$ & $x$ $\vbar$ $\lambda x.e$ $\vbar$ $e$ $e$ & $\lambda$-calculus  \\
			            &     & $\vbar$       & $e\synlink e$                             & linked expression   \\
			            &     & $\vbar$       & $\varepsilon$                             & empty module        \\
			            &     & $\vbar$       & $x\:\texttt{=}\:e$ \texttt{;} $e$         & (recursive) binding
		\end{tabular}
	\end{figure}
\end{frame}
\begin{frame}[c,fragile]
	\frametitle{의미공간 소개}
	\begin{figure}
		\small
		\begin{tabular}{rrcll}
			Environment & $\ctx$ & $\in$         & $\Ctx$                                                                        \\
			Location    & $\ell$ & $\in$         & $\Loc$                                                                        \\
			Value       & $v$    & $\in$         & $\Value \triangleq\Ctx+\ExprVar\times\Expr\times\Ctx$                         \\
			Weak Value  & $w$    & $\in$         & $\Walue\triangleq\Value+\underline\Value$                                     \\
			Environment & $\ctx$ & $\rightarrow$ & $\bullet$                                             & empty stack           \\
			            &        & $\vbar$       & $(x,w)\cons \ctx$                                     & weak value binding    \\
			            &        & $\vbar$       & $(x,\ell)\cons\ctx$                                   & free location binding \\
			Value       & $v$    & $\rightarrow$ & $\ctx$                                                & exported environment  \\
			            &        & $\vbar$       & $\langle \lambda x.e, \ctx \rangle$                   & closure               \\
			Weak Value  & $w$    & $\rightarrow$ & $v$                                                   & value                 \\
			            &        & $\vbar$       & $\mu\ell.v$                                           & recursive value
		\end{tabular}
	\end{figure}
\end{frame}
\begin{frame}[c,fragile]
	\frametitle{의미공간 확장}
	\begin{figure}
		\small
		\begin{tabular}{rrcll}
			Event       & $\event$ & $\rightarrow$ & $\InitE$           & initial environment \\
			            &          & $\vbar$       & $\ReadE(\event,x)$ & read event          \\
			            &          & $\vbar$       & $\CallE(\event,v)$ & call event          \\
			Environment & $\ctx$   & $\rightarrow$ & $\cdots$                                 \\
			            &          & $\vbar$       & $[\event]$         & answer to an event  \\
			Value       & $v$      & $\rightarrow$ & $\cdots$                                 \\
			            &          & $\vbar$       & $\event$           & answer to an event
		\end{tabular}
	\end{figure}
\end{frame}
\begin{frame}[c,fragile]
	\frametitle{실행의미 소개}
	\begin{figure}
		\footnotesize
		\begin{flushright}
			\fbox{$\evjudg{\ctx}{e}{v}$}
		\end{flushright}
		\centering
		\[
			\begin{InfRule}{Id}
				\hypo{\ctx(x)=v}
				\infer1{\evjudg{\ctx}{x}{v}}
			\end{InfRule}\:
			\begin{InfRule}{RecId}
				\hypo{\ctx(x)=\mu\ell.v}
				\infer1{\evjudg{\ctx}{x}{v[\mu\ell.v/\ell]}}
			\end{InfRule}\:
			\begin{InfRule}{Fn}
				\infer0{\evjudg{\ctx}{\lambda x.e}{\langle\lambda x.e, \ctx\rangle}}
			\end{InfRule}\:
			\begin{InfRule}{App}
				\hypo{\evjudg{\ctx}{e_1}{\langle\lambda x.e, \ctx_1\rangle}}
				\hypo{\evjudg{\ctx}{e_2}{v_2}}
				\hypo{\evjudg{(x, v_2)\cons \ctx_1}{e}{v}}
				\infer{2,1}{\evjudg{\ctx}{e_1\:e_2}{v}}
			\end{InfRule}
		\]

		\[
			\begin{InfRule}{Link}
				\hypo{\evjudg{\ctx}{e_1}{\ctx_1}}
				\hypo{\evjudg{\ctx_1}{e_2}{v}}
				\infer2{\evjudg{\ctx}{{e_1}\synlink{e_2}}{v}}
			\end{InfRule}\:
			\begin{InfRule}{Empty}
				\infer0{\evjudg{\ctx}{\varepsilon}{\bullet}}
			\end{InfRule}\:
			\begin{InfRule}{Bind}
				\hypo{\ell\not\in\FLoc(\ctx)}
				\hypo{\evjudg{(x,\ell)\cons\ctx}{e_1}{v_1}}
				\hypo{\evjudg{(x, \mu\ell.v_1)\cons \ctx}{e_1}{\ctx_2}}
				\infer{2,1}{\evjudg{\ctx}{x\:\texttt{=}\:e_1; e_2}{(x,\mu\ell.v_1)\cons\ctx_2}}
			\end{InfRule}
		\]
	\end{figure}
\end{frame}
\begin{frame}
	\frametitle{실행의미 확장}
	\begin{figure}
		\footnotesize
		\[
			\begin{InfRule}{LinkEvent}
				\hypo{\evjudg{\ctx}{e_1}{\event}}
				\hypo{\evjudg{[\event]}{e_2}{v}}
				\infer2{\evjudg{\ctx}{e_1\synlink e_2}{v}}
			\end{InfRule}\:
			\begin{InfRule}{AppEvent}
				\hypo{\evjudg{\ctx}{e_1}{\event}}
				\hypo{\evjudg{\ctx}{e_2}{v}}
				\infer2{\evjudg{\ctx}{e_1\:e_2}{\CallE(\event,v)}}
			\end{InfRule}
		\]

		\[
			\begin{InfRule}{BindEvent}
				\hypo{\ell\not\in\FLoc(\ctx)}
				\hypo{\evjudg{(x,\ell)\cons\ctx}{e_1}{v_1}}
				\hypo{\evjudg{(x, \mu\ell.v_1)\cons \ctx}{e_1}{\event_2}}
				\infer{2,1}{\evjudg{\ctx}{x\:\texttt{=}\:e_1; e_2}{(x,\mu\ell.v_1)\cons[\event_2]}}
			\end{InfRule}
		\]
	\end{figure}
\end{frame}
\begin{frame}
	\frametitle{합치기}
	\footnotesize
	\begin{align*}
		\fbox{$\ctx_0\semlink\cdot\in\Event\rightarrow\pset{\Value}$}                                                                                                                        \\
		\ctx_0\semlink\InitE\triangleq           & \: \{\ctx_0\}                                                                                                                             \\
		\ctx_0\semlink\ReadE(\event,x)\triangleq & \: \{v_+|\ctx_+\in\ctx_0\semlink\event\land\ctx_+(x)=v_+\}                                                                                \\
		\cup                                     & \: \{v_+[\mu\ell.v_+/\ell]|\ctx_+\in\ctx_0\semlink E\land\ctx_+(x)=\mu\ell.v_+\}                                                          \\
		\ctx_0\semlink\CallE(\event,v)\triangleq & \: \{v_+'|\langle\lambda x.e,\ctx_+\rangle\in\ctx_0\semlink E\land v_+\in\ctx_0\semlink v\land(x,v_+)\cons\ctx_+\vdash e\Downarrow v_+'\} \\
		\cup                                     & \: \{\CallE(\event_+,v_+)|\event_+\in\ctx_0\semlink E\land v_+\in\ctx_0\semlink v\}                                                       \\
		\fbox{$\ctx_0\semlink\cdot\in\Ctx\rightarrow\pset{\Ctx}$}                                                                                                                            \\
		\cdots                                                                                                                                                                               \\
		\fbox{$\ctx_0\semlink\cdot\in\Value\rightarrow\pset{\Value}$}                                                                                                                        \\
		\cdots                                                                                                                                                                               \\
		\fbox{$\ctx_0\semlink\cdot\in\Walue\rightarrow\pset{\Walue}$}                                                                                                                        \\
		\ctx_0\semlink\mu\ell.v\triangleq        & \: \{\mu\ell'.v_+|\ell'\not\in\FLoc(v)\cup\FLoc(\ctx_0)\land v_+\in\ctx_0\semlink v[\ell'/\ell]\}
	\end{align*}
\end{frame}
\begin{frame}
	\frametitle{Advance}
	{\small
		\[\eval(e,\ctx)\triangleq\{v|\ctx\vdash e\Downarrow v\}\quad\eval(e,\Sigma)\triangleq\bigcup_{\ctx\in\Sigma}\eval(e,\ctx)\quad\Sigma_0\semlink W\triangleq\bigcup_{\substack{\ctx_0\in\Sigma_0\\w\in W}}(\ctx_0\semlink w)\]
	}
	\begin{thm}[Advance]
		\[\eval(e,\Sigma_0\semlink\Sigma)\subseteq\Sigma_0\semlink\eval(e,\Sigma)\]
	\end{thm}
	\begin{itemize}
		\item 증명은 $\evjudg{\ctx}{e}{v}$에 대한 귀납법으로
		\item Coq으로 엄검증 완료! $\checkmark$
	\end{itemize}
\end{frame}
\begin{frame}[c,fragile]
	\frametitle{따로분석}
	\begin{enumerate}
		\item 증가하는 $\gamma\in\pset{\Walue}\rightarrow\Abs\Walue$
		\item 안전한 $\Abs\eval$: $\Sigma_0\subseteq\gamma(\Abs\ctx_0)\Rightarrow\eval(e,\Sigma_0)\subseteq\gamma(\Abs\eval(e,\Abs\ctx_0))$
		\item 안전한 $\Abs\semlink$: $\Sigma_0\subseteq\gamma(\Abs\ctx_0)\text{ 이고 }W\subseteq\gamma(\Abs{w})\Rightarrow\Sigma_0\semlink W\subseteq\gamma(\Abs\ctx_0\Abs\semlink\Abs{w})$
	\end{enumerate}
	\[\Sigma_0\subseteq\gamma(\Abs\ctx_0)\text{ 이고 }\Sigma\subseteq\gamma(\Abs\ctx)\Rightarrow\eval(e,\Sigma_0\semlink\Sigma)\subseteq\gamma(\Abs\ctx_0\Abs\semlink\Abs\eval(e,\Abs\ctx))\]
\end{frame}
\begin{frame}[c,fragile]
	\frametitle{특히}
	\[\Sigma_0\subseteq\gamma(\Abs\ctx_0)\]
	이고
	\[[\InitE]\in\gamma(\Abs\InitE)\]
	이면
	\[\eval(e_1\synlink e_2,\Sigma_0)\subseteq\gamma(\underbracket{\Abs\eval(e_1,\Abs\ctx_0)}_{\text{따로}}\Abs\semlink\underbracket{\Abs\eval(e_2,\Abs\InitE)}_{\text{따로}})\]
\end{frame}
\begin{frame}[c,fragile]
	\frametitle{예시: CFA}
	\begin{figure}
		\small
		\begin{tabular}{rrcl}
			Program point        & $p$                            & $\in$         & $\mathbb{P}$                                                                    \\
			Abstract event       & $\Abs{\event}$                 & $\in$         & $\Abs{\Event}$                                                                  \\
			Abstract environment & $\Abs{\ctx}$                   & $\in$         & $\Abs{\Ctx}\triangleq(\fin{\ExprVar}{\pset{\mathbb{P}}})\times2^{\Abs{\Event}}$ \\
			Abstract closure     & $\langle\lambda x.p,p'\rangle$ & $\in$         & $\Abs{\text{Clos}}\triangleq\ExprVar\times\mathbb{P}\times\mathbb{P}$           \\
			Abstract value       & $\Abs{v}$                      & $\in$         & $\Abs{\Value}\triangleq\Abs{\Ctx}\times\pset{\Abs{\text{Clos}}}$                \\
			Abstract semantics   & $\Abs\ptree$                   & $\in$         & $\Abs{\mathbb{T}}\triangleq\mathbb{P}\rightarrow\Abs{\Ctx}\times\Abs{\Value}$   \\
			Abstract event       & $\Abs{\event}$                 & $\rightarrow$ & $\Abs{\InitE}$ | $\Abs{\ReadE}(p,x)$ | $\Abs{\CallE}(p,p)$
		\end{tabular}
	\end{figure}
\end{frame}
\begin{frame}[c,fragile]
	\frametitle{직관}
	\[\Abs\ptree=p\mapsto(\Abs\ctx,\Abs{v})\]
	모든 $p$마다, 입력 $\Abs\ctx$와 출력 $\Abs{v}$가 달려있음.
	\[\Abs{v}=(\underbracket{(x\mapsto\{p\},\overbracket{\{\Abs\InitE,\Abs\ReadE(p,x),\Abs\CallE(p_1,p_2)\}}^{\Abs\event\text{ 부분}})}_{\Abs\ctx\text{ 부분}},\underbracket{\{\langle\lambda x.p,p'\rangle\}}_{\Abs{\text{Clos}}\text{ 부분}})\]
	$x\mapsto\{p\}$: $\ctx(x)$가 $\llabel{p}{\ell}$이거나 $p$의 출력.

	$\Abs\ReadE(p,x)$: $p$에 들어오는 $[E]$에서 $x$를 읽음.

	$\Abs\CallE(p_1,p_2)$: $p_1$에서 나온 $E$를 $p_2$에서 나온 $v$에 적용.

	$\langle\lambda x.p,p'\rangle$: $p'$에 들어온 $\ctx$에서 실행되는 함수.
\end{frame}
\begin{frame}[c,fragile]
	\frametitle{$\gamma$ 정의하기}
	\begin{figure}
		\scriptsize
		\begin{flushright}
			\fbox{$\ctx\preceq(\Abs\ctx,\Abs\ptree)$}
		\end{flushright}
		\[
			\begin{InfRule}{Conc-Nil}
				\infer0{\bullet\preceq\Abs\ctx}
			\end{InfRule}\:
			\begin{InfRule}{Conc-ENil}
				\hypo{\event\preceq(\Abs\ctx,\varnothing)}
				\infer1{[\event]\preceq\Abs\ctx}
			\end{InfRule}\:
			\begin{InfRule}{Conc-ConsLoc}
				\hypo{p\in\Abs\ctx.1(x)}
				\hypo{\ctx\preceq\Abs\ctx}
				\infer2{(x,\llabel{p}{\ell})\cons\ctx\preceq\Abs\ctx}
			\end{InfRule}\:
			\begin{InfRule}{Conc-ConsWVal}
				\hypo{p\in\Abs\ctx.1(x)}
				\hypo{w\preceq\Abs\ptree(p).2}
				\hypo{\ctx\preceq\Abs\ctx}
				\infer3{(x,w)\cons\ctx\preceq\Abs\ctx}
			\end{InfRule}
		\]
		\begin{flushright}
			\fbox{$w\preceq(\Abs{v},\Abs\ptree)$}
		\end{flushright}
		\[
			\begin{InfRule}{Conc-Clos}
				\hypo{\langle\lambda x.p,p'\rangle\in\Abs{v}.2}
				\hypo{\ctx\preceq\Abs\ptree(p').1}
				\infer2{\langle\lambda x.p,\ctx\rangle\preceq\Abs{v}}
			\end{InfRule}\:
			\begin{InfRule}{Conc-Rec}
				\hypo{v\preceq\Abs\ptree(p).2}
				\hypo{v\preceq\Abs{v}}
				\infer2{\mu\llabel{p}{\ell}.v\preceq\Abs{v}}
			\end{InfRule}\:
			\begin{InfRule}{Conc-Init}
				\hypo{\Abs\InitE\in\Abs{v}.1.2}
				\infer1{\InitE\preceq\Abs{v}}
			\end{InfRule}
		\]

		\[
			\begin{InfRule}{Conc-Read}
				\hypo{\Abs\ReadE(p,x)\in\Abs{v}.1.2}
				\hypo{[\event]\preceq\Abs\ptree(p).1}
				\infer2{\ReadE(\event,x)\preceq\Abs{v}}
			\end{InfRule}\:
			\begin{InfRule}{Conc-Call}
				\hypo{\Abs\CallE(p_1,p_2)\in\Abs{v}.1.2}
				\hypo{\event\preceq\Abs\ptree(p_1).2}
				\hypo{v\preceq\Abs\ptree(p_2).2}
				\infer3{\CallE(\event,v)\preceq\Abs{v}}
			\end{InfRule}
		\]
	\end{figure}
\end{frame}
\begin{frame}
	\frametitle{$\gamma$ 정의하기}
	\[\gamma(\Abs{v},\Abs\ptree)\triangleq\{w|w\preceq(\Abs{v},\Abs\ptree)\}\]
\end{frame}
\begin{frame}[c,fragile]
	\frametitle{안전한 $\Abs\Step$}
	\begin{itemize}
		\item $\Abs\ptree\sqsupseteq\Abs\Step(\Abs\ptree)$이면
		      \[\evjudg{\ctx}{p}{v}\Rightarrow\ctx\in\gamma(\Abs\ptree(p).1,\Abs\ptree)\Rightarrow v\in\gamma(\Abs\ptree(p).2,\Abs\ptree)\]
		      가 성립하는
	\end{itemize}
\end{frame}
\begin{frame}[c,fragile]
	\frametitle{안전한 $\Abs\Link$}
	\begin{itemize}
		\item $\ctx_0\in\gamma(\Abs\ctx_0,\Abs\ptree_0)$인 $\Abs\ctx_0,\Abs\ptree_0$와
		\item $\Abs\ptree$가 주어졌을 때,
		\item $\Abs\ptree_+\sqsupseteq\Abs\ptree_0$이고 $\Abs\ptree_+\sqsupseteq\Abs\Link(\Abs\ctx_0,\Abs\ptree,\Abs\ptree_+)$이면
	\end{itemize}
	\[w_+\in\ctx_0\semlink w\Rightarrow w\in\gamma(\Abs\ptree(p).1,\Abs\ptree)\Rightarrow w_+\in\gamma(\Abs\ptree_+(p).1,\Abs\ptree_+)\]
	와
	\[w_+\in\ctx_0\semlink w\Rightarrow w\in\gamma(\Abs\ptree(p).2,\Abs\ptree)\Rightarrow w_+\in\gamma(\Abs\ptree_+(p).2,\Abs\ptree_+)\]
	가 성립하는
\end{frame}
\begin{frame}
	\frametitle{안전한 $\Abs{\sembracket{p_0}}$와 $\Abs\semlink$}
	\[\Abs{\sembracket{p_0}}(\Abs\ctx_0,\Abs\ptree_0)\triangleq\lfp(\lambda\Abs\ptree.\Abs\Step(\Abs\ptree)\sqcup[p_0\mapsto(\Abs\ctx_0,\bot)]\sqcup\Abs\ptree_0)\]
	\[(\Abs\ctx_0,\Abs\ptree_0)\Abs\semlink\Abs\ptree\triangleq\lfp(\lambda\Abs\ptree_+.\Abs\Link(\Abs\ctx_0,\Abs\ptree,\Abs\ptree_+)\sqcup\Abs\ptree_0)\]
\end{frame}
\begin{frame}
	\begin{center}\LARGE
		감사합니다
	\end{center}
\end{frame}
\end{document}
